%% Generated by Sphinx.
\def\sphinxdocclass{report}
\documentclass[a4paper,11pt,oneside,english,openany]{sphinxmanual}
\ifdefined\pdfpxdimen
   \let\sphinxpxdimen\pdfpxdimen\else\newdimen\sphinxpxdimen
\fi \sphinxpxdimen=.75bp\relax
\ifdefined\pdfimageresolution
    \pdfimageresolution= \numexpr \dimexpr1in\relax/\sphinxpxdimen\relax
\fi
%% let collapsible pdf bookmarks panel have high depth per default
\PassOptionsToPackage{bookmarksdepth=5}{hyperref}

\PassOptionsToPackage{booktabs}{sphinx}
\PassOptionsToPackage{colorrows}{sphinx}

\PassOptionsToPackage{warn}{textcomp}
\usepackage[utf8]{inputenc}
\ifdefined\DeclareUnicodeCharacter
% support both utf8 and utf8x syntaxes
  \ifdefined\DeclareUnicodeCharacterAsOptional
    \def\sphinxDUC#1{\DeclareUnicodeCharacter{"#1}}
  \else
    \let\sphinxDUC\DeclareUnicodeCharacter
  \fi
  \sphinxDUC{00A0}{\nobreakspace}
  \sphinxDUC{2500}{\sphinxunichar{2500}}
  \sphinxDUC{2502}{\sphinxunichar{2502}}
  \sphinxDUC{2514}{\sphinxunichar{2514}}
  \sphinxDUC{251C}{\sphinxunichar{251C}}
  \sphinxDUC{2572}{\textbackslash}
\fi
\usepackage{cmap}
\usepackage[T1]{fontenc}
\usepackage{amsmath,amssymb,amstext}
\usepackage{babel}





\usepackage[Bjarne]{fncychap}
\usepackage{sphinx}

\fvset{fontsize=auto}
\usepackage{geometry}


% Include hyperref last.
\usepackage{hyperref}
% Fix anchor placement for figures with captions.
\usepackage{hypcap}% it must be loaded after hyperref.
% Set up styles of URL: it should be placed after hyperref.
\urlstyle{same}

\addto\captionsenglish{\renewcommand{\contentsname}{Contents:}}

\usepackage{sphinxmessages}
\setcounter{tocdepth}{1}


		\renewcommand{\familydefault}{\rmdefault}
        \usepackage{setspace} % Required for more spacing options
		\parindent 0px
		\setstretch{1.1}
        \usepackage[T1]{fontenc}
        \usepackage{charter}
        \usepackage{amsmath}
        \usepackage{amsfonts}
        \usepackage{amssymb}
        \usepackage{fancyhdr}
        \usepackage{arydshln}
        \usepackage{xcolor}
        \usepackage{makeidx}  % To create an index
        \usepackage{tocloft}  % To customize Table of Contents

        % Fancy header and footer
        \pagestyle{fancy}
        \fancyhead[L]{\leftmark}
        \fancyhead[R]{\thepage}
        \fancyfoot[C]{}

        % Table of Contents: Formatting (removing subsections)
        \setcounter{tocdepth}{1}  % Limit ToC to sections only (remove subsections)
        \renewcommand{\cftsecfont}{\normalfont\bfseries}  % Bold section titles
        \renewcommand{\cftsecindent}{0em}                 % No indentation for sections
        \setlength{\cftbeforesecskip}{0.5ex}              % Space before section titles

        % Index Configuration (simplified to sections only)
        \makeindex  % Generate index

        % Page layout adjustments
		\usepackage{geometry}
        \geometry{top=1.5cm, bottom=1.5cm, left=1.5cm, right=1.5cm}
        \usepackage{hyperref}
        \hypersetup{colorlinks=true,linkcolor=violet,filecolor=magenta,urlcolor=blue}

        % Better verbatim environment for code blocks
        %\usepackage{fvextra}
        \usepackage{tcolorbox}
        
        \newtcbox{\tbox}{colframe=black,colback=gray!5,boxrule=0.7pt,arc=4pt,boxsep=5pt}
    

\title{Smooth Component Filtering (SCF)}
\date{May 21, 2026}
\release{1.0.0}
\author{AE,SB,BB}
\newcommand{\sphinxlogo}{\vbox{}}
\renewcommand{\releasename}{Release}
\makeindex
\begin{document}

\ifdefined\shorthandoff
  \ifnum\catcode`\=\string=\active\shorthandoff{=}\fi
  \ifnum\catcode`\"=\active\shorthandoff{"}\fi
\fi

\pagestyle{empty}
\sphinxmaketitle
\pagestyle{plain}
\sphinxtableofcontents
\pagestyle{normal}
\phantomsection\label{\detokenize{index::doc}}
\sphinxstepscope



\sphinxstepscope


\chapter{scf module}
\label{\detokenize{scf:module-scf}}\label{\detokenize{scf:scf-module}}\label{\detokenize{scf::doc}}\index{module@\spxentry{module}!scf@\spxentry{scf}}\index{scf@\spxentry{scf}!module@\spxentry{module}}

\section{Smooth Component Filtering (SCF)}
\label{\detokenize{scf:smooth-component-filtering-scf}}
\sphinxAtStartPar
By \sphinxstylestrong{convolved gridded visibility} \(\mathcal{V}_{cg}(\nu_n)\) and given a \sphinxstylestrong{suitable window function} \(H(n)\) of \sphinxstylestrong{specified smoothing scale} (controlled by parameter \(N\)), we calculate the following \sphinxhyphen{}
\begin{enumerate}
\sphinxsetlistlabels{\arabic}{enumi}{enumii}{}{.}%
\item {} 
\sphinxAtStartPar
\(\left(\mathcal{V}_{cg}\ast H\right)(\nu_n) = \displaystyle \sum_m\, \mathcal{V}_{cg}(\nu_m)H(n-m)\) and,

\item {} 
\sphinxAtStartPar
\(\left(\rm U\ast H\right)(\nu_n)    = \displaystyle\sum_m\, {\rm U}(\nu_m)H(n-m)\)

\end{enumerate}

\sphinxAtStartPar
Dividing them we get :
\begin{equation*}
\begin{split}\begin{align}
    \boxed{\mathcal{V}_{cg}^S(\nu_n) = \left(\mathcal{V}_{cg}\ast H\right)(\nu_n)/\left(\rm U\ast H\right)(\nu_n)}
\end{align}\end{split}
\end{equation*}
\sphinxAtStartPar
We refer \(\mathcal{V}_{cg}^S(\nu_n)\) as \sphinxstylestrong{Smooth Component} of the convolved gridded visibility \(\mathcal{V}_{cg}(\nu_n)\) and \sphinxstylestrong{subtracting} it from the \(\mathcal{V}_{cg}(\nu_n)\) gives us the \sphinxstylestrong{filtered component}.
\begin{equation*}
\begin{split}\begin{align}
    \boxed{\mathcal{V}_{cg}^F(\nu_n) = \mathcal{V}_{cg}(\nu_n) - \mathcal{V}_{cg}^S(\nu_n)}
\end{align}\end{split}
\end{equation*}
\sphinxAtStartPar
\(\mathcal{V}_{cg}(\nu_i) = 0\) for the channels \(\nu_i\), that are \sphinxstylestrong{flagged} and the divison by \(\left(\rm U\ast H\right)(\nu_n)\) is necessary to \sphinxstylestrong{ensure correct normalization after smoothing}, i.e. \(\left(\mathcal{V}_{cg}\ast H\right)(\nu_n)\). Here \({\rm U}(\nu_i) = 0\) for the flagged channels and \({\rm U}(\nu_i) = 1\) otherwise. The \sphinxstylestrong{smoothing is restricted to a smaller range near the boundaries} and to account for this \(N\) no of \sphinxstylestrong{channels are discarded} from both ends of \(\mathcal{V}_{cg}^F(\nu_n)\) to account this.

\sphinxAtStartPar
If \(\Delta\nu_c\) denotes the \sphinxstylestrong{channel separation} then for a given \(N\) corresponds to  \sphinxstylestrong{smoothing scale} of \(N\Delta\nu_c\). For example \(N = 25, 50\, \text{and}\, 75\) and \(\Delta\nu_c = 0.04 \rm MHz\) corresponds to smoothing scales of \(N\Delta\nu_c = 1, 2\, \text{and}\, 3\, \rm MHz\) respectively.
\subsubsection*{Notes}

\sphinxAtStartPar
We expect \sphinxstylestrong{SCF} to \sphinxstylestrong{suppress} the values of \(|P(k_\perp,k_\parallel)|\) somewhere between the half and full width of the filter, i.e.
\begin{equation*}
\begin{split}\left[k_{\parallel}\right]_{\rm F}/2 \leq k_\parallel \leq \left[k_{\parallel}\right]_{\rm F}\end{split}
\end{equation*}
\sphinxAtStartPar
Where \(\left[k_{\parallel}\right]_{\rm F} = \dfrac{2\pi}{r'N\Delta\nu_c}\).


\section{Window functions}
\label{\detokenize{scf:window-functions}}
\sphinxAtStartPar
This module supports \sphinxstylestrong{four types of window functions} namely \sphinxhyphen{}
\begin{enumerate}
\sphinxsetlistlabels{\arabic}{enumi}{enumii}{}{.}%
\item {} 
\sphinxAtStartPar
\sphinxstylestrong{Hann/Hanning} window function \(H_{\rm Hann}(n)\).

\item {} 
\sphinxAtStartPar
\sphinxstylestrong{Hamming} window function \(H_{\rm Hamming}(n)\).

\item {} 
\sphinxAtStartPar
\sphinxstylestrong{Blackman} window function \(H_{\rm Blackman}(n)\).

\item {} 
\sphinxAtStartPar
\sphinxstylestrong{Kaiser} window function \(H_{\rm Kaiser}(n)\) with \(\beta = 14\) .

\end{enumerate}

\sphinxAtStartPar
The functional form’s are given as :
\begin{equation*}
\begin{split}\setlength{\fboxrule}{0.7pt}
\setlength{\fboxsep}{12pt}     % EXTRA padding
\fcolorbox{black}{gray!3}{$
\begin{aligned}
    H_{\rm Hann}(n) &= \frac{1}{2N}\left[0.5+0.5\cos\left(2\pi\frac{n}{2N}\right)\right] \quad& -N \leq n \leq N\\[1em]
    H_{\rm Hamming}(n) &= \frac{1}{2N}\left[0.54+0.46\cos\left(2\pi\frac{n}{2N}\right)\right] \quad& -N \leq n \leq N\\[1em]
    H_{\rm Blackman}(n) &= \frac{1}{2N}\left[0.42+0.5\cos\left(2\pi\frac{n}{2N}\right)+0.08\cos\left(4\pi\frac{n}{2N}\right)\right] \quad& -N \leq n \leq N\\[1em]
    H_{\rm Kaiser}(n) &= \frac{1}{2N}\frac{I_0\left(\sqrt{1-\left(\frac{n}{N}\right)^2}\right)}{I_0(\beta)} \quad& -N \leq n \leq N
\end{aligned}$}\end{split}
\end{equation*}
\begin{figure}[H]
\centering
\includegraphics[width=1.05\linewidth]{../../../windows_123.pdf}
\caption{Different window functions with different smoothing scales.}
\end{figure}

\clearpage


\section{Module functions:}
\label{\detokenize{scf:module-functions}}\index{doscf() (in module scf)@\spxentry{doscf()}\spxextra{in module scf}}

\begin{fulllineitems}
\phantomsection\label{\detokenize{scf:scf.doscf}}
\pysigstartsignatures
\pysiglinewithargsret{\sphinxcode{\sphinxupquote{scf.}}\sphinxbfcode{\sphinxupquote{doscf}}}{\sphinxparam{\DUrole{n}{GV}}\sphinxparamcomma \sphinxparam{\DUrole{n}{SM}}\sphinxparamcomma \sphinxparam{\DUrole{n}{window}\DUrole{o}{=}\DUrole{default_value}{\textquotesingle{}hann\textquotesingle{}}}\sphinxparamcomma \sphinxparam{\DUrole{n}{method}\DUrole{o}{=}\DUrole{default_value}{\textquotesingle{}auto\textquotesingle{}}}}{}
\pysigstopsignatures
\sphinxAtStartPar
Given gridded visibility array it performs \sphinxstylestrong{SCF (Smooth Component Filtering)} given a \sphinxstylestrong{window function} and \sphinxstylestrong{suitable smoothing scale}. It \sphinxstylestrong{subtracts out} the \sphinxstylestrong{smooth part} of it and returns \sphinxstylestrong{residual filtered visibilities}.
\begin{quote}\begin{description}
\sphinxlineitem{Parameters}\begin{itemize}
\item {} 
\sphinxAtStartPar
\sphinxstyleliteralstrong{\sphinxupquote{GV}} (\sphinxstyleliteralemphasis{\sphinxupquote{np.ndarray}}) \textendash{} Convolved\sphinxhyphen{}gridded visibility array. The last dimension must have shape \sphinxstylestrong{(,768)}.

\item {} 
\sphinxAtStartPar
\sphinxstyleliteralstrong{\sphinxupquote{SM}} (\sphinxstyleliteralemphasis{\sphinxupquote{int}}\sphinxstyleliteralemphasis{\sphinxupquote{ or }}\sphinxstyleliteralemphasis{\sphinxupquote{float}}) \textendash{} \sphinxstylestrong{Smoothing scale} in MHz of specified window function.

\item {} 
\sphinxAtStartPar
\sphinxstyleliteralstrong{\sphinxupquote{window}} (\sphinxstyleliteralemphasis{\sphinxupquote{str \{\textquotesingle{}hann\textquotesingle{}}}\sphinxstyleliteralemphasis{\sphinxupquote{,}}\sphinxstyleliteralemphasis{\sphinxupquote{\textquotesingle{}hamming\textquotesingle{}}}\sphinxstyleliteralemphasis{\sphinxupquote{,}}\sphinxstyleliteralemphasis{\sphinxupquote{\textquotesingle{}blackman\textquotesingle{}}}\sphinxstyleliteralemphasis{\sphinxupquote{,}}\sphinxstyleliteralemphasis{\sphinxupquote{\textquotesingle{}kaiser\textquotesingle{}\}}}\sphinxstyleliteralemphasis{\sphinxupquote{, }}\sphinxstyleliteralemphasis{\sphinxupquote{optional}}) \textendash{} Window function name. Valid ones are \sphinxstylestrong{‘hann’, ‘hamming’, ‘blackman’, ‘kaiser’}. If not any of this it kills the current iteration. By default ‘hann’ window.

\item {} 
\sphinxAtStartPar
\sphinxstyleliteralstrong{\sphinxupquote{method}} (\sphinxstyleliteralemphasis{\sphinxupquote{str \{\textquotesingle{}auto\textquotesingle{}}}\sphinxstyleliteralemphasis{\sphinxupquote{,}}\sphinxstyleliteralemphasis{\sphinxupquote{\textquotesingle{}direct\textquotesingle{}}}\sphinxstyleliteralemphasis{\sphinxupquote{,}}\sphinxstyleliteralemphasis{\sphinxupquote{\textquotesingle{}fft\textquotesingle{}\}}}\sphinxstyleliteralemphasis{\sphinxupquote{, }}\sphinxstyleliteralemphasis{\sphinxupquote{optional}}) \textendash{} The \sphinxstylestrong{method} by which the \sphinxstylestrong{convolution} is performed. Either by \sphinxstylestrong{‘direct’ or ‘fft’}. By \sphinxstylestrong{default} is \sphinxstylestrong{‘auto’}, means \sphinxstylestrong{whichever is faster is applied}.

\end{itemize}

\sphinxlineitem{Returns}
\sphinxAtStartPar
\sphinxstylestrong{GV\_filtered} \textendash{} The \sphinxstylestrong{filtered visibility} which is obtained by \sphinxstylestrong{subtracting out the smooth part}. Shape same as input GV array \sphinxstylestrong{except in the last dimension} in which first and last \({\rm NW} = {\rm SM}/{\Delta\nu_c}\) no of channels are sliced. Where \(\Delta\nu_c\) is the separation between two frequency channels in MHz. By default the value is \sphinxstyleemphasis{0.04}.

\sphinxlineitem{Return type}
\sphinxAtStartPar
np.ndarray

\end{description}\end{quote}
\subsubsection*{Examples}

\begin{sphinxVerbatim}[commandchars=\\\{\}]
\PYG{g+gp}{\PYGZgt{}\PYGZgt{}\PYGZgt{} }\PYG{k+kn}{import} \PYG{n+nn}{numpy} \PYG{k}{as} \PYG{n+nn}{np}
\PYG{g+gp}{\PYGZgt{}\PYGZgt{}\PYGZgt{} }\PYG{k+kn}{import} \PYG{n+nn}{scf}
\PYG{g+go}{Imported SCF, Import Time : 21/05/26 | 03:00:01 PM}
\PYG{g+gp}{\PYGZgt{}\PYGZgt{}\PYGZgt{} }\PYG{n}{GV} \PYG{o}{=} \PYG{n}{np}\PYG{o}{.}\PYG{n}{load}\PYG{p}{(}\PYG{l+s+s2}{\PYGZdq{}}\PYG{l+s+s2}{/home/baijayanta/GV\PYGZus{}7200\PYGZus{}pool.npy}\PYG{l+s+s2}{\PYGZdq{}}\PYG{p}{)} \PYG{c+c1}{\PYGZsh{} GV array}
\PYG{g+gp}{\PYGZgt{}\PYGZgt{}\PYGZgt{} }\PYG{n+nb}{print}\PYG{p}{(}\PYG{l+s+sa}{f}\PYG{l+s+s2}{\PYGZdq{}}\PYG{l+s+si}{\PYGZob{}}\PYG{n}{GV}\PYG{o}{.}\PYG{n}{shape}\PYG{+w}{ }\PYG{l+s+si}{= \PYGZcb{}}\PYG{l+s+s2}{\PYGZdq{}}\PYG{p}{)}
\PYG{g+go}{GV.shape = (2, 46438, 768)}
\PYG{g+gp}{\PYGZgt{}\PYGZgt{}\PYGZgt{} }\PYG{n}{SM} \PYG{o}{=} \PYG{l+m+mi}{2} \PYG{c+c1}{\PYGZsh{} smoothing scale in MHz}
\PYG{g+gp}{\PYGZgt{}\PYGZgt{}\PYGZgt{} }\PYG{n}{NW} \PYG{o}{=} \PYG{n+nb}{int}\PYG{p}{(}\PYG{n}{SM}\PYG{o}{/}\PYG{l+m+mf}{0.04}\PYG{p}{)}
\PYG{g+gp}{\PYGZgt{}\PYGZgt{}\PYGZgt{} }\PYG{n+nb}{print}\PYG{p}{(}\PYG{l+s+sa}{f}\PYG{l+s+s2}{\PYGZdq{}}\PYG{l+s+si}{\PYGZob{}}\PYG{n}{NW}\PYG{+w}{ }\PYG{l+s+si}{= \PYGZcb{}}\PYG{l+s+s2}{\PYGZdq{}}\PYG{p}{)}
\PYG{g+go}{NW = 50}
\PYG{g+gp}{\PYGZgt{}\PYGZgt{}\PYGZgt{} }\PYG{n}{NW} \PYG{o}{=} \PYG{l+m+mi}{50}
\PYG{g+gp}{\PYGZgt{}\PYGZgt{}\PYGZgt{} }\PYG{c+c1}{\PYGZsh{} perform SCF}
\PYG{g+gp}{\PYGZgt{}\PYGZgt{}\PYGZgt{} }\PYG{n}{GV\PYGZus{}filtered} \PYG{o}{=} \PYG{n}{scf}\PYG{o}{.}\PYG{n}{doscf}\PYG{p}{(}\PYG{n}{GV}\PYG{p}{,} \PYG{n}{SM}\PYG{p}{,} \PYG{n}{window} \PYG{o}{=} \PYG{l+s+s1}{\PYGZsq{}}\PYG{l+s+s1}{hann}\PYG{l+s+s1}{\PYGZsq{}}\PYG{p}{,} \PYG{n}{method} \PYG{o}{=} \PYG{l+s+s1}{\PYGZsq{}}\PYG{l+s+s1}{fft}\PYG{l+s+s1}{\PYGZsq{}}\PYG{p}{)}
\PYG{g+go}{\PYGZlt{}============= Performing SCF =============\PYGZgt{}}
\PYG{g+go}{\PYGZlt{}=== Start Time: 21/05/26 | 03:00:16 PM ===\PYGZgt{}}
\PYG{g+go}{Convoluting kernel : fft}
\PYG{g+go}{Window function    : Hanning}
\PYG{g+go}{Smoothing scale    : 2 MHz}
\PYG{g+go}{\PYGZlt{}=== End Time  : 21/05/26 | 03:00:18 PM ===\PYGZgt{}}
\PYG{g+go}{\PYGZlt{}=== Elapsed   : 00:00:02.53, 2.533 seconds.}
\PYG{g+go}{\PYGZlt{}================ Done SCF ================\PYGZgt{}}
\PYG{g+gp}{\PYGZgt{}\PYGZgt{}\PYGZgt{} }\PYG{n+nb}{print}\PYG{p}{(}\PYG{l+s+sa}{f}\PYG{l+s+s2}{\PYGZdq{}}\PYG{l+s+si}{\PYGZob{}}\PYG{n}{GV\PYGZus{}filtered}\PYG{o}{.}\PYG{n}{shape}\PYG{+w}{ }\PYG{l+s+si}{= \PYGZcb{}}\PYG{l+s+s2}{\PYGZdq{}}\PYG{p}{)} \PYG{c+c1}{\PYGZsh{} 768\PYGZhy{}2*NW}
\PYG{g+go}{GV\PYGZus{}filtered.shape = (2, 46438, 668)}
\end{sphinxVerbatim}

\begin{figure}[H]
\centering
\includegraphics[width=1.05\linewidth]{../../../Flag.pdf}
\caption{The Flagging of MWA data. The black colored regions are flagged.}
\end{figure}

\begin{figure}[H]
\centering
\includegraphics[width=1.05\linewidth]{../../../scf.pdf}
\caption{Real part of $\mathcal{V}_{cg},\mathcal{V}_{cg}^{S},\mathcal{V}_{cg}^{F}$ as a function of frequency $\nu$ for a fixed grid point corresponding to $\ell_g = 261$. This is shown only for XX polarization.}
\end{figure}

\end{fulllineitems}




\renewcommand{\indexname}{Index}
\printindex
\end{document}